\documentclass[11pt]{article}
\usepackage[a4paper,margin=2.3cm]{geometry}
\usepackage{amsmath,amssymb,amsfonts,amsthm,mathtools}
\usepackage{bm}
\usepackage[most]{tcolorbox}
\usepackage{hyperref}
\hypersetup{colorlinks=true,linkcolor=blue!55!black,urlcolor=blue!55!black}
\allowdisplaybreaks[4]

\newcommand{\mf}{\mathfrak}
\newcommand{\eps}{\varepsilon}
\newcommand{\dg}{\dagger}
\newcommand{\ip}[1]{\int\!\frac{d^{3}#1}{(2\pi)^{3}}}
\newcommand{\dthree}{\delta^{(3)}}
\newcommand{\ket}[1]{\left|#1\right\rangle}
\newcommand{\bra}[1]{\left\langle#1\right|}
\newcommand{\vac}{\ket{0}}
\newcommand{\OM}{\Omega}
\newcommand{\half}{\tfrac12}
\newcommand{\acomm}[2]{\big\{#1,\,#2\big\}}

\newtcolorbox{res}[1][]{colback=blue!3,colframe=blue!45!black,boxrule=0.6pt,
  arc=2pt,left=6pt,right=6pt,top=4pt,bottom=4pt,#1}
\newtcolorbox{warn}[1][]{colback=orange!5,colframe=orange!60!black,boxrule=0.6pt,
  arc=2pt,left=6pt,right=6pt,top=4pt,bottom=4pt,#1}

\title{\bf Coefficients of the quark--gluon evolution operator $\OM_{q\bar q g}$\\
fixed by unitarity through $O(g^{2})$}
\date{}
\author{}

\begin{document}
\maketitle
\vspace{-1.3cm}

\begin{abstract}
\noindent
Starting from the normal ordered ansatz for $\OM_{q\bar qg}$ and from the coefficients
$\mf A_1$, $\mf A_3$, $\mf B_1$, $\mf B_3$, $\mf C_3$, $\mf C_4$, $\mf D_1$, $\mf D_3$ determined by light-cone
perturbation theory, we impose $\OM^{\dg}\OM=1$ order by order in $g$ and solve for the remaining
coefficients. We obtain $\mf C_0$, $\mf A_2$, $\mf A_4$, $\mf B_2$, $\mf B_4$, $\mf C_1$,
$\mf C_2$ and $\mf D_2$ in closed form. Two by-products of the same algebra are recorded: a
consistency condition that the LCPT result for $\mf D_3$ must satisfy (only the anti-Hermitian
part of $\mf D_3$ is independent), and the appearance of an $a^{\dg}a$ structure --- the quark
loop --- which is absent from the ansatz and which must be added to the Yang--Mills coefficient
$B_3$ of the pure gauge paper. We show that its diagonal part is proportional to a scaleless
transverse integral and therefore vanishes in dimensional regularization. The relation for
$\mf A_2$ is verified independently against the explicit LCPT wave functions, including the exact
cancellation of the instantaneous contributions.
\end{abstract}

\tableofcontents
\bigskip

%==================================================================
\section{Setup and conventions}
%==================================================================

\subsection{The two sectors of $\OM$}
\label{sec:twosectors}

The full evolution operator acts on the soft Hilbert space built from soft gluons, quarks and
antiquarks. It contains a pure gauge part and the quark part,
\begin{equation}
\OM=\OM_{\rm YM}+\OM_{q\bar qg}-1 ,
\label{eq:split}
\end{equation}
the subtraction avoiding a double counting of the identity ($\OM_{\rm YM}$ carries the
normalization $N$ and $\OM_{q\bar qg}$ carries $\mf C_0$).

\begin{warn}
\textbf{This point matters for every result below.} Unitarity is a statement about the
\emph{full} $\OM$, not about $\OM_{q\bar qg}$ separately. Consequently the quark coefficients
fixed by unitarity are \emph{not} determined by quark coefficients alone: they mix with the two
$O(g)$ coefficients of the pure gauge sector, the Weizs\"acker--Williams coefficient $A^b_j(p)$
and the three-gluon coefficient $C^{dcb}_{lkj}(p,q,r)$. Every relation for
$\mf A_2,\mf B_2,\mf B_4,\mf C_1,\mf C_2$ below contains such cross terms.

If one instead insisted that $\OM_{q\bar qg}$ be unitary by itself (i.e.\ $\OM=\OM_{\rm YM}\cdot
\OM_{q}$ with each factor unitary) all these cross terms would disappear and one would get the much
simpler $\mf A_2=-\mf A_1^{\dg}$. That factorized convention is \emph{inconsistent} with defining
the coefficients by direct matching to the LCWF components, as is done here and in the pure gauge
paper: Appendix~\ref{app:check} verifies by explicit LCPT computation that the LCWF-matched
$\mf A_1$ and $\mf A_2$ obey the relation \emph{with} the cross term, and not $\mf A_2=-\mf
A_1^\dg$.
\end{warn}

We write, using the notation of the pure gauge paper for the $O(g)$ gluonic coefficients,
\begin{equation}
\OM=1+\OM_{1}+\OM_{2}+O(g^{3}),\qquad \OM_1=O(g),\quad \OM_2=O(g^{2}),
\end{equation}
\begin{align}
\OM_{1}=&\ip{p}\Big[A^{b}_{j}(p)\,a^{\dg b}_{j}(p)-A^{\dg b}_{j}(p)\,a^{b}_{j}(p)\Big]
\nonumber\\
&+\ip{p}\ip{q}\ip{r}\Big[C^{dcb}_{lkj}(p,q,r)\,a^{\dg d}_{l}(r)a^{\dg c}_{k}(q)a^{b}_{j}(p)
-C^{\dg bcd}_{jkl}(r,q,p)\,a^{\dg d}_{l}(r)a^{c}_{k}(q)a^{b}_{j}(p)\Big]
\nonumber\\
&+\ip{p}\ip{l}\ip{m}\Big[\mf A^{b\alpha\beta}_{3j\lambda_{2}\lambda_{1}}(p,m,l)\,
b^{\beta\dg}_{\lambda_{1}}(l)d^{\alpha\dg}_{\lambda_{2}}(m)a^{b}_{j}(p)
\nonumber\\&\hspace{3.1cm}
+\mf A^{\beta\alpha b}_{4\lambda_{1}\lambda_{2}j}(l,m,p)\,
a^{\dg b}_{j}(p)d^{\alpha}_{\lambda_{2}}(m)b^{\beta}_{\lambda_{1}}(l)\Big].
\label{eq:Om1}
\end{align}
Here $A^b_j(p)$ and $C^{dcb}_{lkj}(p,q,r)$ are respectively the coefficients called $A_2$ and
$C_1$ in the pure gauge paper, and we have already used the pure gauge $O(g)$ unitarity relations
$A_1=-A^\dg_2$ and $C^{dcb}_{2lkj}(p,q,r)=-C^{\dg bcd}_{1jkl}(r,q,p)$. All momentum integrals are
$\int d^{3}p/(2\pi)^{3}$ and all repeated indices are summed. We abbreviate
\begin{equation}
\delta_{pq}^{bc,jk}\equiv(2\pi)^{3}\delta^{bc}\delta_{jk}\dthree(p-q),
\end{equation}
which is what a contraction $[a^{b}_{j}(p),a^{\dg c}_{k}(q)]$ produces.

\subsection{Index and argument ordering}
\label{sec:convention}

Every coefficient in the ansatz is written with its colour superscripts, its
polarization/helicity subscripts and its momentum arguments listed \emph{in the order in which the
corresponding operators appear when the operator string is read from right to left}. For instance
\begin{equation}
\mf B^{cb\alpha\beta}_{1kj\lambda_{2}\lambda_{1}}(q,p,m,l)\;\;
b^{\beta\dg}_{\lambda_{1}}(l)\,d^{\alpha\dg}_{\lambda_{2}}(m)\,
a^{b}_{j}(p)\,a^{c}_{k}(q),
\end{equation}
since reading the operator string from right to left gives
$(c,k,q)$, $(b,j,p)$, $(\alpha,\lambda_2,m)$, $(\beta,\lambda_1,l)$.
We verified that this rule is obeyed by \emph{all} of the LCPT-determined coefficients
$\mf A_1$, $\mf A_3$, $\mf B_1$, $\mf B_3$, $\mf C_1$, $\mf C_2$, $\mf C_3$, $\mf D_1$, $\mf D_3$ exactly as printed.
Four coefficients in the draft ansatz do not obey it, and since these are precisely the ones we
are solving for, we fix them here:

\begin{center}
\begin{tabular}{lll}
\hline
coefficient & as printed & used here\\
\hline
$\mf A_4$ & $\mf A^{\beta\alpha b}_{4\lambda_1\lambda_2 j}(p)$ &
$\mf A^{\beta\alpha b}_{4\lambda_1\lambda_2 j}(l,m,p)$ (arguments restored)\\
$\mf B_2$ & $\mf B^{\beta\alpha cb}_{2\lambda_1\lambda_2kj}(l,m,p,q)$ &
$\mf B^{\beta\alpha cb}_{2\lambda_1\lambda_2kj}(l,m,q,p)$\\
$\mf B_4$ & $\mf B^{cb\beta\alpha}_{4kj\lambda_1\lambda_2}(m,l,p,q)$ &
$\mf B^{cb\beta\alpha}_{4kj\lambda_1\lambda_2}(q,p,l,m)$\\
$\mf C_4$ & $\mf C^{bcde\beta\alpha}_{4jklm\lambda_1\lambda_2}(p,q,r,s,l,m)$ &
$\mf C^{bcde\alpha\beta}_{4jkon\lambda_2\lambda_1}(p,q,r,s,m,l)$ (as in $\mf C_3$)\\
\hline
\end{tabular}
\end{center}

Two further notational points. First, the draft uses the symbol $m$ both for the antiquark
momentum and for the polarization index of $a^{\dg e}_{m}(s)$ in the $\mf C$ terms. We rename the
gluon polarization indices in the $\mf C$ sector to
\begin{equation}
a^{\dg e}_{n}(s)\,a^{\dg d}_{o}(r)\,a^{c}_{k}(q)\,a^{b}_{j}(p),
\end{equation}
i.e.\ $(j,k,o,n)$ for the momenta $(p,q,r,s)$, reserving $l,m$ for the quark and antiquark momenta.

\begin{warn}
\textbf{Second, and more substantively: $\mf D_2$ as printed cannot close under unitarity.} The
printed structure is
\[
\mf D_{2}\;a^{\dg b}_{j}(p)\,a^{c}_{k}(q)\,d^{\alpha'}_{\lambda'_{2}}(m')b^{\beta'}_{\lambda'_{1}}(l')
d^{\alpha}_{\lambda_{2}}(m)b^{\beta}_{\lambda_{1}}(l),
\]
with one gluon created and one annihilated. The Hermitian conjugate of the $\mf D_1$ structure
$b^{\dg}d^{\dg}b^{\dg}d^{\dg}\,a(p)a(q)$ is $a^{\dg}(q)a^{\dg}(p)\,d'b'db$, with \emph{two}
gluons created. Moreover the printed structure does not conserve any quantum number at $O(g^2)$:
annihilating two $q\bar q$ pairs while destroying one gluon and creating one is a $g^{4}$ process.
We therefore read the second operator as $a^{\dg c}_{k}(q)$,
\[
\mf D_{2}\;a^{\dg b}_{j}(p)\,a^{\dg c}_{k}(q)\,d^{\alpha'}_{\lambda'_{2}}(m')b^{\beta'}_{\lambda'_{1}}(l')
d^{\alpha}_{\lambda_{2}}(m)b^{\beta}_{\lambda_{1}}(l),
\]
which is the correct $O(g^{2})$ partner of $\mf D_{1}$ (two $q\bar q$ pairs annihilating into two
gluons). All index and argument labels are then consistent with the right-to-left rule exactly as
printed.
\end{warn}

\subsection{Power counting}

As stated: $\mf A_{3},\mf A_{4}=O(g)$; all of
$\mf A_{1},\mf A_{2},\mf B_{1..4},\mf C_{1..4},\mf D_{1,2,3}$ are $O(g^{2})$; and
$\mf C_{0}=1+O(g^{2})$. On the pure gauge side $A^b_j$ and $C^{dcb}_{lkj}$ are $O(g)$.

%==================================================================
\section{The unitarity hierarchy}
%==================================================================

Inserting $\OM=1+\OM_{1}+\OM_{2}+\dots$ into $\OM^{\dg}\OM=1$ and collecting powers of $g$,
\begin{equation}
\OM^{\dg}\OM=1+\big(\OM_{1}+\OM^{\dg}_{1}\big)
+\big(\OM_{2}+\OM^{\dg}_{2}+\OM^{\dg}_{1}\OM_{1}\big)+O(g^{3})=1,
\end{equation}
so that
\begin{equation}
\boxed{\;\OM^{\dg}_{1}=-\OM_{1}\;},
\qquad
\boxed{\;\OM_{2}+\OM^{\dg}_{2}=-\OM^{\dg}_{1}\OM_{1}=\OM_{1}^{2}\;}.
\label{eq:hierarchy}
\end{equation}
The second equality uses the first. Note that Eq.~\eqref{eq:hierarchy} determines only the
\emph{Hermitian part} of $\OM_{2}$; the anti-Hermitian part is the information supplied by LCPT.
For a coefficient whose operator structure and the structure of its Hermitian conjugate are
\emph{different} (e.g.\ $\mf B_{1}$ and $\mf B_{2}$), this splitting is equivalent to: know one,
get the other. For a coefficient that is its own conjugate partner (e.g.\ $\mf D_{3}$, or the
identity $\mf C_{0}$) unitarity only constrains half of it, and becomes a consistency check on the
LCPT result. We use both facts below.

Because the coefficients are functionals of the non-Abelian valence charge density
$\rho^{a}$, obeying $[\rho^{a}(k),\rho^{b}(p)]=if^{abc}\rho^{c}(k+p)$, they do \emph{not} commute
with one another. The order of the coefficient factors on the right-hand sides below is
significant and has been tracked throughout: in $\OM_1^2=\OM_1\OM_1$ the coefficient coming from
the \emph{left} factor stands to the left. Soft operators ($a,b,d$ and their conjugates) commute
with all valence operators.

\subsection{Enumeration of the contractions}
\label{sec:enum}

Write $\OM_{1}=G+Q$ with $G$ the pure gauge part (first two lines of Eq.~\eqref{eq:Om1}) and $Q$
the quark part (third line), both anti-Hermitian. Then
\begin{equation}
\OM_{1}^{2}=\underbrace{G^{2}}_{\text{pure gauge}}+\;GQ+QG+Q^{2}.
\end{equation}
$G^{2}$ reproduces the relations of the pure gauge paper and is dropped. Normal ordering
$GQ+QG+Q^{2}$ and using the elementary contractions
\begin{equation}
a^{b}_{j}(p)a^{\dg c}_{k}(q)=\delta^{bc,jk}_{pq}+a^{\dg c}_{k}(q)a^{b}_{j}(p),
\end{equation}
\begin{equation}
d^{\alpha}_{\lambda_{2}}(m)b^{\beta}_{\lambda_{1}}(l)\,
b^{\beta'\dg}_{\lambda'_{1}}(l')d^{\alpha'\dg}_{\lambda'_{2}}(m')
=\delta_{\kappa\kappa'}\delta_{\bar\kappa\bar\kappa'}
-\delta_{\kappa\kappa'}d^{\alpha'\dg}_{\lambda'_{2}}(m')d^{\alpha}_{\lambda_{2}}(m)
-\delta_{\bar\kappa\bar\kappa'}b^{\beta'\dg}_{\lambda'_{1}}(l')b^{\beta}_{\lambda_{1}}(l)
+b^{\beta'\dg}b^{\beta}d^{\alpha'\dg}d^{\alpha},
\label{eq:fermicontr}
\end{equation}
(with $\kappa=\{l,\lambda_{1},\beta\}$, $\bar\kappa=\{m,\lambda_{2},\alpha\}$) gives the following
exhaustive list of terms. We label them (a)--(t) and indicate the normal ordered structure each
produces.

\begin{center}
\renewcommand{\arraystretch}{1.15}
\begin{tabular}{clll}
\hline
& product & structure(s) produced & feeds\\
\hline
(a)&$(\mf A_{3}b^{\dg}d^{\dg}a)(\mf A_{3}b^{\dg}d^{\dg}a)$ & $b^{\dg}d^{\dg}b^{\dg}d^{\dg}aa$ & $\mf D_{1}$\\
(b)&$(\mf A_{3}b^{\dg}d^{\dg}a)(-\mf A^{\dg}_{3}a^{\dg}db)$ & $b^{\dg}d^{\dg}bd$; $b^{\dg}d^{\dg}a^{\dg}a\,db$ & $\mf D_{3}$; $\times$\\
(c)&$(-\mf A^{\dg}_{3}a^{\dg}db)(\mf A_{3}b^{\dg}d^{\dg}a)$ & $a^{\dg}a$; $a^{\dg}a\,b^{\dg}b$; $a^{\dg}a\,d^{\dg}d$ & $\times$\\
(d)&$(-\mf A^{\dg}_{3}a^{\dg}db)(-\mf A^{\dg}_{3}a^{\dg}db)$ & $a^{\dg}a^{\dg}\,dbdb$ & $\mf D_{2}$\\
(e)&$(Aa^{\dg})(\mf A_{3}b^{\dg}d^{\dg}a)$ & $a^{\dg}a\,b^{\dg}d^{\dg}$ & $\mf B_{3}$\\
(f)&$(\mf A_{3}b^{\dg}d^{\dg}a)(Aa^{\dg})$ & $b^{\dg}d^{\dg}$; $a^{\dg}a\,b^{\dg}d^{\dg}$ & $\mf A_{1}$; $\mf B_{3}$\\
(g)&$(-A^{\dg}a)(\mf A_{3}b^{\dg}d^{\dg}a)$ & $b^{\dg}d^{\dg}aa$ & $\mf B_{1}$\\
(h)&$(\mf A_{3}b^{\dg}d^{\dg}a)(-A^{\dg}a)$ & $b^{\dg}d^{\dg}aa$ & $\mf B_{1}$\\
(i)&$(Ca^{\dg}a^{\dg}a)(\mf A_{3}b^{\dg}d^{\dg}a)$ & $b^{\dg}d^{\dg}a^{\dg}a^{\dg}aa$ & $\mf C_{4}$\\
(j)&$(\mf A_{3}b^{\dg}d^{\dg}a)(Ca^{\dg}a^{\dg}a)$ & $a^{\dg}a\,b^{\dg}d^{\dg}$; $\mf C_{4}$-structure & $\mf B_{3}$; $\mf C_{4}$\\
(k)&$(-C^{\dg}a^{\dg}aa)(\mf A_{3}b^{\dg}d^{\dg}a)$ & $b^{\dg}d^{\dg}a^{\dg}aaa$ & $\mf C_{3}$\\
(l)&$(\mf A_{3}b^{\dg}d^{\dg}a)(-C^{\dg}a^{\dg}aa)$ & $b^{\dg}d^{\dg}aa$; $\mf C_{3}$-structure & $\mf B_{1}$; $\mf C_{3}$\\
(m)&$(-\mf A^{\dg}_{3}a^{\dg}db)(Aa^{\dg})$ & $a^{\dg}a^{\dg}db$ & $\mf B_{2}$\\
(n)&$(Aa^{\dg})(-\mf A^{\dg}_{3}a^{\dg}db)$ & $a^{\dg}a^{\dg}db$ & $\mf B_{2}$\\
(o)&$(-\mf A^{\dg}_{3}a^{\dg}db)(-A^{\dg}a)$ & $a^{\dg}a\,db$ & $\mf B_{4}$\\
(p)&$(-A^{\dg}a)(-\mf A^{\dg}_{3}a^{\dg}db)$ & $db$; $a^{\dg}a\,db$ & $\mf A_{2}$; $\mf B_{4}$\\
(q)&$(-\mf A^{\dg}_{3}a^{\dg}db)(Ca^{\dg}a^{\dg}a)$ & $a^{\dg}a^{\dg}a^{\dg}a\,db$ & $\mf C_{1}$\\
(r)&$(Ca^{\dg}a^{\dg}a)(-\mf A^{\dg}_{3}a^{\dg}db)$ & $a^{\dg}a^{\dg}db$; $\mf C_{1}$-structure & $\mf B_{2}$; $\mf C_{1}$\\
(s)&$(-\mf A^{\dg}_{3}a^{\dg}db)(-C^{\dg}a^{\dg}aa)$ & $a^{\dg}a^{\dg}aa\,db$ & $\mf C_{2}$\\
(t)&$(-C^{\dg}a^{\dg}aa)(-\mf A^{\dg}_{3}a^{\dg}db)$ & $a^{\dg}a\,db$; $\mf C_{2}$-structure & $\mf B_{4}$; $\mf C_{2}$\\
\hline
\end{tabular}
\end{center}
Entries marked $\times$ produce operator structures that are \emph{not present} in the ansatz;
they are discussed in Section~\ref{sec:missing}.

Note that no term produces the identity: every quark structure in $\OM_{1}$ carries at least one
soft fermion operator, and the only fully contracted products leave behind either $a^{\dg}a$
(from (c)) or $b^{\dg}d^{\dg}bd$ (from (b)). This is the origin of $\mf C_{0}=1$.

%==================================================================
\section{Results}
%==================================================================

Throughout, each relation is preceded by the normal ordered operator string that both sides
multiply, so that no ordering ambiguity remains.

\subsection{Order $g$: $\mf A_{4}$}

From $\OM^{\dg}_{1}=-\OM_{1}$, using
$\big(b^{\beta\dg}_{\lambda_{1}}(l)d^{\alpha\dg}_{\lambda_{2}}(m)a^{b}_{j}(p)\big)^{\dg}
=a^{\dg b}_{j}(p)d^{\alpha}_{\lambda_{2}}(m)b^{\beta}_{\lambda_{1}}(l)$:

\begin{res}
\emph{Structure} $a^{\dg b}_{j}(p)\,d^{\alpha}_{\lambda_{2}}(m)\,b^{\beta}_{\lambda_{1}}(l)$:
\begin{equation}
\boxed{\;
\mf A^{\beta\alpha b}_{4\lambda_{1}\lambda_{2}j}(l,m,p)
=-\Big[\mf A^{b\alpha\beta}_{3j\lambda_{2}\lambda_{1}}(p,m,l)\Big]^{\dg}
\;}
\label{eq:A4}
\end{equation}
\end{res}

\noindent
This is the only $O(g)$ relation. Note that it pairs $\mf A_{4}$ with $\mf A_{3}$, and (below)
$\mf A_{2}$ with $\mf A_{1}$. The pairings $(\mf A_{3},\mf A_{1})$ and $(\mf A_{4},\mf A_{2})$
quoted in the earlier notes cannot hold: $\mf A_{1},\mf A_{2}$ carry no gluon operator while
$\mf A_{3},\mf A_{4}$ carry one, so they cannot be related by Hermitian conjugation. They are
also of different orders in $g$.

\subsection{$\mf C_{0}$}

As shown at the end of Section~\ref{sec:enum}, the identity component of $\OM_{1}^{2}$ receives no
quark contribution. Hence

\begin{res}
\begin{equation}
\boxed{\;\mf C_{0}=1+O(g^{4})\;}
\label{eq:C0}
\end{equation}
The first correction appears at $O(g^{4})$, where the $O(g^{4})$ unitarity condition
$\OM_{4}+\OM^{\dg}_{4}=-\big(\OM^{\dg}_{1}\OM_{3}+\OM^{\dg}_{2}\OM_{2}+\OM^{\dg}_{3}\OM_{1}\big)$
gives, from the identity projection of $\OM^{\dg}_{2}\OM_{2}$,
\begin{equation}
\mf C_{0}=1-\half\ip{l}\ip{m}
\Big[\mf A^{\alpha\beta}_{1\lambda_{2}\lambda_{1}}(m,l)\Big]^{\dg}
\mf A^{\alpha\beta}_{1\lambda_{2}\lambda_{1}}(m,l)+O(g^{6}).
\end{equation}
\end{res}

\noindent
The physical statement is the expected one: $\OM\vac=\mf C_{0}\vac+\int \mf A_{1}\ket{q\bar q}+\dots$,
and since pair creation out of the valence background first occurs at $O(g^{2})$, the
depletion of the vacuum component is $O(g^{4})$.

\subsection{$\mf A_{2}$}

The only contribution is (p). Explicitly,
\begin{align}
(\text{p})&=\ip{p}\ip{q}A^{\dg b}_{j}(p)\Big[\mf A^{c\alpha\beta}_{3k\lambda_{2}\lambda_{1}}(q,m,l)\Big]^{\dg}
a^{b}_{j}(p)a^{\dg c}_{k}(q)\,d^{\alpha}_{\lambda_{2}}(m)b^{\beta}_{\lambda_{1}}(l)
\nonumber\\
&\supset\ip{p}A^{\dg b}_{j}(p)\Big[\mf A^{b\alpha\beta}_{3j\lambda_{2}\lambda_{1}}(p,m,l)\Big]^{\dg}
d^{\alpha}_{\lambda_{2}}(m)b^{\beta}_{\lambda_{1}}(l).
\end{align}
On the left-hand side, $[\OM_{2}]_{db}=\mf A_{2}$ and
$[\OM^{\dg}_{2}]_{db}=[\mf A_{1}]^{\dg}$, since
$\big(\mf A_{1}b^{\dg}_{\lambda_1}(l)d^{\dg}_{\lambda_2}(m)\big)^{\dg}
=d_{\lambda_2}(m)b_{\lambda_1}(l)\,\mf A^{\dg}_{1}$.

\begin{res}
\emph{Structure} $d^{\alpha}_{\lambda_{2}}(m)\,b^{\beta}_{\lambda_{1}}(l)$:
\begin{equation}
\boxed{\;
\mf A^{\beta\alpha}_{2\lambda_{1}\lambda_{2}}(l,m)
=-\Big[\mf A^{\alpha\beta}_{1\lambda_{2}\lambda_{1}}(m,l)\Big]^{\dg}
+\ip{p}A^{\dg b}_{j}(p)\,\Big[\mf A^{b\alpha\beta}_{3j\lambda_{2}\lambda_{1}}(p,m,l)\Big]^{\dg}
\;}
\label{eq:A2}
\end{equation}
Equivalently, taking the Hermitian conjugate,
\begin{equation}
\mf A^{\alpha\beta}_{1\lambda_{2}\lambda_{1}}(m,l)
+\Big[\mf A^{\beta\alpha}_{2\lambda_{1}\lambda_{2}}(l,m)\Big]^{\dg}
=\ip{p}\mf A^{b\alpha\beta}_{3j\lambda_{2}\lambda_{1}}(p,m,l)\,A^{b}_{j}(p).
\label{eq:A2conj}
\end{equation}
\end{res}

\noindent
The second form has a transparent reading: the amplitude for creating a $q\bar q$ pair out of the
vacuum, plus the conjugate of the amplitude for absorbing one, equals the amplitude for the
valence charge to emit a gluon ($A$) which then splits ($\mf A_{3}$). The ordering is fixed:
$\mf A_{3}$ stands to the \emph{left} of $A$, because the gluon is emitted before it splits.
Equation~\eqref{eq:A2conj} is verified by explicit LCPT computation in
Appendix~\ref{app:check}, including the non-trivial cancellation of the instantaneous
$\rho\cdot\mathcal{J}_q$ contributions between $\mf A_1$ and $\mf A_2^\dg$.

\subsection{$\mf B_{2}$}

Contributions (m), (n) and the contracted part of (r). For (r),
\begin{align}
(\text{r})\Big|_{\delta}
&=-\ip{u}C^{bcb'}_{jkj'}(u,q,p)\Big[\mf A^{b'\alpha\beta}_{3j'\lambda_{2}\lambda_{1}}(u,m,l)\Big]^{\dg}
\;a^{\dg b}_{j}(p)a^{\dg c}_{k}(q)\,d^{\alpha}_{\lambda_{2}}(m)b^{\beta}_{\lambda_{1}}(l).
\end{align}
The left-hand side is $\mf B_{2}+[\mf B_{1}]^{\dg}$, since
$\big(\mf B_{1}b^{\dg}(l)d^{\dg}(m)a^{b}_{j}(p)a^{c}_{k}(q)\big)^{\dg}
=a^{\dg c}_{k}(q)a^{\dg b}_{j}(p)d(m)b(l)\,\mf B_{1}^{\dg}$.

\begin{res}
\emph{Structure} $a^{\dg b}_{j}(p)\,a^{\dg c}_{k}(q)\,d^{\alpha}_{\lambda_{2}}(m)\,b^{\beta}_{\lambda_{1}}(l)$:
\begin{align}
\boxed{
\begin{aligned}
\mf B^{\beta\alpha cb}_{2\lambda_{1}\lambda_{2}kj}(l,m,q,p)
=&-\Big[\mf B^{cb\alpha\beta}_{1kj\lambda_{2}\lambda_{1}}(q,p,m,l)\Big]^{\dg}
\\
&-\Big[\mf A^{b\alpha\beta}_{3j\lambda_{2}\lambda_{1}}(p,m,l)\Big]^{\dg}A^{c}_{k}(q)
-A^{b}_{j}(p)\Big[\mf A^{c\alpha\beta}_{3k\lambda_{2}\lambda_{1}}(q,m,l)\Big]^{\dg}
\\
&-\ip{u}C^{bcb'}_{jkj'}(u,q,p)\Big[\mf A^{b'\alpha\beta}_{3j'\lambda_{2}\lambda_{1}}(u,m,l)\Big]^{\dg}
\end{aligned}}
\label{eq:B2}
\end{align}
\end{res}

\noindent
Both sides are automatically symmetric under the simultaneous exchange
$(b,j,p)\leftrightarrow(c,k,q)$, as they must be since $a^{\dg}a^{\dg}$ commute: the two $A\mf
A^{\dg}_{3}$ terms map into each other, and $C^{dcb}_{lkj}(p,q,r)$ is symmetric in its two
creation labels $(d,l,r)\leftrightarrow(c,k,q)$.

\subsection{$\mf B_{4}$}

Contributions (o), the uncontracted part of (p), and the contracted part of (t). In (t) the
$C^{\dg}$ term carries two annihilation operators, either of which can contract with the $a^{\dg}$
of $\mf A^{\dg}_{3}$; since $C^{\dg bcd}_{jkl}(r,q,p)$ is symmetric under
$(b,j,p)\leftrightarrow(c,k,q)$ the two contractions are equal and give a factor $2$.

\begin{res}
\emph{Structure} $a^{\dg b}_{j}(p)\,a^{c}_{k}(q)\,d^{\alpha}_{\lambda_{2}}(m)\,b^{\beta}_{\lambda_{1}}(l)$:
\begin{align}
\boxed{
\begin{aligned}
\mf B^{cb\beta\alpha}_{4kj\lambda_{1}\lambda_{2}}(q,p,l,m)
=&-\Big[\mf B^{\alpha\beta bc}_{3\lambda_{2}\lambda_{1}jk}(m,l,p,q)\Big]^{\dg}
\\
&+\Big[\mf A^{b\alpha\beta}_{3j\lambda_{2}\lambda_{1}}(p,m,l)\Big]^{\dg}A^{\dg c}_{k}(q)
+A^{\dg c}_{k}(q)\Big[\mf A^{b\alpha\beta}_{3j\lambda_{2}\lambda_{1}}(p,m,l)\Big]^{\dg}
\\
&+2\ip{u}C^{\dg b'cb}_{j'kj}(p,q,u)\Big[\mf A^{b'\alpha\beta}_{3j'\lambda_{2}\lambda_{1}}(u,m,l)\Big]^{\dg}
\end{aligned}}
\label{eq:B4}
\end{align}
\end{res}

\noindent
The first two terms on the second line form the anticommutator
$\acomm{[\mf A_{3}]^{\dg}}{A^{\dg}}$ of two valence operators; it is \emph{not} $2A^{\dg}\mf
A^{\dg}_{3}$, because $[\rho,\rho]\neq0$.

\subsection{$\mf C_{1}$}

Contributions (q) and the uncontracted part of (r). Both have the structure ``$\mf A^{\dg}_{3}$
carries one of the three creation labels, $C$ carries the other two plus the annihilation label'',
differing only in the order of the two valence operators. Let $\mathcal S_{3}$ denote
symmetrization over the three creation labels $(c,k,q)$, $(d,o,r)$, $(e,n,s)$ (the coefficient of
a product of three commuting $a^{\dg}$'s is defined as its symmetric part).

\begin{res}
\emph{Structure} $a^{\dg e}_{n}(s)\,a^{\dg d}_{o}(r)\,a^{\dg c}_{k}(q)\,a^{b}_{j}(p)\,
d^{\alpha}_{\lambda_{2}}(m)\,b^{\beta}_{\lambda_{1}}(l)$:
\begin{align}
\boxed{
\begin{aligned}
\mf C^{\beta\alpha bcde}_{1\lambda_{1}\lambda_{2}jkon}(l,m,p,q,r,s)
=&-\Big[\mf C^{bcde\alpha\beta}_{3jkon\lambda_{2}\lambda_{1}}(p,q,r,s,m,l)\Big]^{\dg}
\\
&-\mathcal S_{3}\,\acomm{\Big[\mf A^{e\alpha\beta}_{3n\lambda_{2}\lambda_{1}}(s,m,l)\Big]^{\dg}}
{C^{dcb}_{okj}(p,q,r)}
\end{aligned}}
\label{eq:C1}
\end{align}
\end{res}

\subsection{$\mf C_{2}$}

Contributions (s) and the uncontracted part of (t), again differing only by the order of the two
valence factors. Here $\mathcal S_{2}$ symmetrizes over the two creation labels $(d,o,r)$,
$(e,n,s)$.

\begin{res}
\emph{Structure} $a^{\dg e}_{n}(s)\,a^{\dg d}_{o}(r)\,a^{c}_{k}(q)\,a^{b}_{j}(p)\,
d^{\alpha}_{\lambda_{2}}(m)\,b^{\beta}_{\lambda_{1}}(l)$:
\begin{align}
\boxed{
\begin{aligned}
\mf C^{\beta\alpha bcde}_{2\lambda_{1}\lambda_{2}jkon}(l,m,p,q,r,s)
=&-\Big[\mf C^{bcde\alpha\beta}_{4jkon\lambda_{2}\lambda_{1}}(p,q,r,s,m,l)\Big]^{\dg}
\\
&+\mathcal S_{2}\,\acomm{\Big[\mf A^{e\alpha\beta}_{3n\lambda_{2}\lambda_{1}}(s,m,l)\Big]^{\dg}}
{C^{\dg bcd}_{jko}(r,q,p)}
\end{aligned}}
\label{eq:C2}
\end{align}
\end{res}

\noindent
The sign difference between Eqs.~\eqref{eq:C1} and \eqref{eq:C2} traces back to the relative sign
between the $C$ and $C^{\dg}$ terms in the anti-Hermitian $\OM_{1}$, Eq.~\eqref{eq:Om1}.

\subsection{$\mf D_{2}$}

The only contribution is (d). Since each $db$ pair contains an even number of fermion operators,
$d(m)b(l)$ and $d(m')b(l')$ commute and no relative sign arises:
\begin{equation}
(\text{d})=\ip{p}\ip{q}
\Big[\mf A^{b\alpha\beta}_{3j\lambda_{2}\lambda_{1}}(p,m,l)\Big]^{\dg}
\Big[\mf A^{c\alpha'\beta'}_{3k\lambda'_{2}\lambda'_{1}}(q,m',l')\Big]^{\dg}
a^{\dg b}_{j}(p)a^{\dg c}_{k}(q)\,d^{\alpha'}_{\lambda'_{2}}(m')b^{\beta'}_{\lambda'_{1}}(l')
d^{\alpha}_{\lambda_{2}}(m)b^{\beta}_{\lambda_{1}}(l).
\end{equation}

\begin{res}
\emph{Structure} $a^{\dg b}_{j}(p)\,a^{\dg c}_{k}(q)\,
d^{\alpha'}_{\lambda'_{2}}(m')b^{\beta'}_{\lambda'_{1}}(l')\,
d^{\alpha}_{\lambda_{2}}(m)b^{\beta}_{\lambda_{1}}(l)$:
\begin{align}
\boxed{
\begin{aligned}
\mf D^{\beta\alpha\beta'\alpha' cb}_{2\lambda_{1}\lambda_{2}\lambda'_{1}\lambda'_{2}kj}(l,m,l',m',q,p)
=&-\Big[\mf D^{cb\alpha'\beta'\alpha\beta}_{1kj\lambda'_{2}\lambda'_{1}\lambda_{2}\lambda_{1}}
(q,p,m',l',m,l)\Big]^{\dg}
\\
&+\mathcal S_{(bjp)(ckq)}\;
\Big[\mf A^{b\alpha\beta}_{3j\lambda_{2}\lambda_{1}}(p,m,l)\Big]^{\dg}
\Big[\mf A^{c\alpha'\beta'}_{3k\lambda'_{2}\lambda'_{1}}(q,m',l')\Big]^{\dg}
\end{aligned}}
\label{eq:D2}
\end{align}
\end{res}

\noindent
The last term is the statement that two $q\bar q$ pairs annihilate into two gluons through two
independent $q\bar q\to g$ transitions; the antisymmetry under
$(\kappa\bar\kappa)\leftrightarrow(\mu\bar\mu)$ required by Fermi statistics is supplied
automatically by the fermionic operators it multiplies.

%==================================================================
\section{Consistency conditions and structures missing from the ansatz}
\label{sec:missing}
%==================================================================

The same algebra produces three further pieces of information. The first is a check on an
LCPT-determined coefficient; the last two concern operator structures absent from the ansatz.

\subsection{$\mf D_{3}$ is only half independent}

$\mf D_{3}$ multiplies $b^{\dg}d^{\dg}b'd'$, and the Hermitian conjugate of that structure is
again of the same type (with the primed and unprimed labels exchanged). Unitarity therefore fixes
its Hermitian part. From the contracted part of (b), using
$d^{\alpha'}_{\lambda'_{2}}(m')b^{\beta'}_{\lambda'_{1}}(l')
=-b^{\beta'}_{\lambda'_{1}}(l')d^{\alpha'}_{\lambda'_{2}}(m')$:

\begin{res}
\emph{Structure} $b^{\beta\dg}_{\lambda_{1}}(l)d^{\alpha\dg}_{\lambda_{2}}(m)\,
b^{\beta'}_{\lambda'_{1}}(l')d^{\alpha'}_{\lambda'_{2}}(m')$:
\begin{align}
\boxed{
\begin{aligned}
&\mf D^{\alpha'\beta'\alpha\beta}_{3\lambda'_{2}\lambda'_{1}\lambda_{2}\lambda_{1}}(m',l',m,l)
+\Big[\mf D^{\alpha\beta\alpha'\beta'}_{3\lambda_{2}\lambda_{1}\lambda'_{2}\lambda'_{1}}(m,l,m',l')\Big]^{\dg}
\\[2pt]
&\hspace{2.2cm}
=\ip{p}\mf A^{b\alpha\beta}_{3j\lambda_{2}\lambda_{1}}(p,m,l)
\Big[\mf A^{b\alpha'\beta'}_{3j\lambda'_{2}\lambda'_{1}}(p,m',l')\Big]^{\dg}
\end{aligned}}
\label{eq:D3}
\end{align}
\end{res}

\noindent
This is a genuine constraint that the LCPT result for $\mf D_{3}$ must satisfy, and it is worth
using as a check: the right-hand side is built entirely from the $O(g)$ splitting amplitude, while
the left-hand side comes from the $q\bar q\to q\bar q$ wave function. Physically it states that
the $s$-channel annihilation graph $q\bar q\to g\to q\bar q$, symmetrized, is fixed by
probability conservation. Only the anti-Hermitian part of $\mf D_{3}$ carries independent LCPT
information.

\subsection{The quark loop: an $a^{\dg}a$ term is missing from the ansatz}

Entry (c) of the table produces, from the fully contracted term in Eq.~\eqref{eq:fermicontr},
\begin{equation}
-\ip{l}\ip{m}\sum_{\lambda_{1}\lambda_{2}\alpha\beta}
\Big[\mf A^{b\alpha\beta}_{3j\lambda_{2}\lambda_{1}}(p,m,l)\Big]^{\dg}
\mf A^{c\alpha\beta}_{3k\lambda_{2}\lambda_{1}}(q,m,l)\;\;a^{\dg b}_{j}(p)a^{c}_{k}(q),
\label{eq:quarkloop}
\end{equation}
a purely gluonic, gluon-number-preserving structure. $\OM_{q\bar qg}$ as written contains no such
term, so Eq.~\eqref{eq:quarkloop} has nowhere to go. The resolution is that it is a quark
contribution to the Yang--Mills coefficient $B_{3}$ of the pure gauge paper (the coefficient of
$a^{\dg c}_{k}(q)a^{b}_{j}(p)$), which in the presence of quarks must be split as
$B_{3}=B_{3}^{\rm YM}+B_{3}^{(q)}$ with
\begin{equation}
B_{3}^{(q)}+\big[B_{3}^{(q)}\big]^{\dg}
=-\ip{l}\ip{m}\Big[\mf A^{b\alpha\beta}_{3j\lambda_{2}\lambda_{1}}(p,m,l)\Big]^{\dg}
\mf A^{c\alpha\beta}_{3k\lambda_{2}\lambda_{1}}(q,m,l).
\end{equation}
This is the familiar statement that the probability for a soft gluon to remain a gluon is reduced
by the probability of its splitting into a $q\bar q$ pair.

Its momentum-diagonal part is the quark contribution to the normalization of the one-gluon state,
and it can be evaluated in closed form. Write the $g\to q\bar q$ vertex structure of the ansatz as
\begin{equation}
\Gamma^{i}(k;l,m)=\frac{2k^{i}}{k^{+}}-\frac{(\bm\sigma\!\cdot\!\bm l)\sigma^{i}}{l^{+}}
-\frac{\sigma^{i}(\bm\sigma\!\cdot\!\bm m)}{m^{+}},\qquad k=l+m,
\end{equation}
and introduce $z=l^{+}/k^{+}$ and the relative transverse momentum
$\bm\kappa=\bm l-z\bm k=(m^{+}\bm l-l^{+}\bm m)/k^{+}$. Then, remarkably, all dependence on the
total momentum $\bm k$ cancels:

\begin{res}
\textbf{Lemma.} With $\bm l=z\bm k+\bm\kappa$, $\bm m=(1-z)\bm k-\bm\kappa$,
\begin{equation}
\boxed{\;
\Gamma^{i}(k;l,m)=\frac{(2z-1)\kappa^{i}+i\,\epsilon^{ij}\kappa^{j}\sigma^{3}}{k^{+}z(1-z)}\;}
\label{eq:lemma}
\end{equation}
\end{res}

\begin{proof}[Proof]
Substituting and using $\{\sigma^{j},\sigma^{i}\}=2\delta^{ij}$,
\begin{equation}
2k^{i}-(\bm\sigma\!\cdot\!\bm k)\sigma^{i}-\sigma^{i}(\bm\sigma\!\cdot\!\bm k)=0,
\end{equation}
so the three terms proportional to $\bm k$ cancel identically and
\begin{equation}
\Gamma^{i}=\frac{1}{k^{+}}\Big[-\frac{(\bm\sigma\!\cdot\!\bm\kappa)\sigma^{i}}{z}
+\frac{\sigma^{i}(\bm\sigma\!\cdot\!\bm\kappa)}{1-z}\Big].
\end{equation}
With $\sigma^{i}\sigma^{j}=\delta^{ij}+i\epsilon^{ij}\sigma^{3}$ one has
$(\bm\sigma\!\cdot\!\bm\kappa)\sigma^{i}=\kappa^{i}-i\epsilon^{ij}\kappa^{j}\sigma^{3}$ and
$\sigma^{i}(\bm\sigma\!\cdot\!\bm\kappa)=\kappa^{i}+i\epsilon^{ij}\kappa^{j}\sigma^{3}$, which
gives Eq.~\eqref{eq:lemma}.
\end{proof}

\noindent
Using Eq.~\eqref{eq:lemma}, $\mathrm{Tr}\,\sigma^{3}=0$ and
$\epsilon^{ij}\epsilon^{i'j'}\kappa^{j}\kappa^{j'}=\delta^{ii'}\kappa^{2}-\kappa^{i}\kappa^{i'}$,
\begin{equation}
\sum_{\lambda_{1}\lambda_{2}}\Phi^{i\,*}_{\lambda_1\lambda_2}\Phi^{i'}_{\lambda_1\lambda_2}
=\mathrm{Tr}\big[\Gamma^{i\dg}\Gamma^{i'}\big]
=\frac{2\big[\delta^{ii'}\kappa^{2}-4z(1-z)\kappa^{i}\kappa^{i'}\big]}{(k^{+})^{2}z^{2}(1-z)^{2}}
\;\xrightarrow[\ \int d^{2}\kappa\ ]{}\;
\frac{2\,\delta^{ii'}\kappa^{2}\big[z^{2}+(1-z)^{2}\big]}{(k^{+})^{2}z^{2}(1-z)^{2}},
\end{equation}
where $\Phi^{i}_{\lambda_1\lambda_2}=\chi^{\dg}_{\lambda_1}\Gamma^{i}\chi_{\lambda_2}$ and we used
$1-2z(1-z)=z^{2}+(1-z)^{2}$. Together with the light-cone energy denominator
\begin{equation}
\eps_{k}-\eps_{l}-\eps_{m}=-\frac{\kappa^{2}}{2z(1-z)k^{+}},
\end{equation}
which enters $\mf A_{3}$ squared, every power of $\kappa$ cancels except an overall $1/\kappa^{2}$:
\begin{equation}
\big[\text{Eq.}~\eqref{eq:quarkloop}\big]_{p=q}
\;\propto\;g^{2}N_{f}T_{F}\,\delta^{bc}\delta_{jk}\dthree(p-q)
\int_{0}^{1}\!dz\,\big[z^{2}+(1-z)^{2}\big]\int\!\frac{d^{2}\kappa}{\kappa^{2}} .
\label{eq:scaleless}
\end{equation}

\begin{res}
The transverse integral in Eq.~\eqref{eq:scaleless} is scaleless and vanishes in dimensional
regularization. \textbf{The quark loop therefore gives no correction to the normalization of the
soft one-gluon state at $O(g^{2})$}, exactly as the gluon loop does not in the pure gauge case.
Note also that, unlike the gluonic $g\to gg$ case, the longitudinal integral
$\int_{0}^{1}dz\,[z^{2}+(1-z)^{2}]=2/3$ is finite: there is no $\log(\vee/\Lambda)$ from the
quark loop.
\end{res}

\subsection{Structures generated but not present in the ansatz}

For completeness, the partial contractions in Eq.~\eqref{eq:fermicontr} also generate
\begin{equation}
a^{\dg}a\,b^{\dg}b,\qquad a^{\dg}a\,d^{\dg}d\qquad\text{(from (c))},
\qquad
b^{\dg}d^{\dg}a^{\dg}a\,db\qquad\text{(from (b))} .
\end{equation}
These are legitimate $O(g^{2})$ structures. They contribute only when $\OM$ acts on a state
containing \emph{both} soft gluons and a soft $q\bar q$ pair, so they are invisible in the
matching sectors used here (vacuum, one, two and three gluons, and $q\bar q$) and can consistently
be omitted --- but they should be restored before $\OM$ is applied to mixed states.

\subsection{Caveat: the $O(g)$ structures $q\to qg$ and $\bar q\to\bar q g$}
\label{sec:FG}

\begin{warn}
The results above are complete \emph{for the ansatz as given}. However, the same vertex that
produces $\mf A_{3}$ ($g\to q\bar q$) also produces, at the same order $O(g)$, the emission of a
gluon by a soft quark or antiquark. The corresponding normal ordered structures
\begin{equation}
\mf F^{b\beta'\beta}_{j\lambda'_{1}\lambda_{1}}(p,l',l)\,a^{\dg b}_{j}(p)b^{\beta'\dg}_{\lambda'_{1}}(l')b^{\beta}_{\lambda_{1}}(l)
\;-\;\mathrm{h.c.},
\qquad
\mf G^{b\alpha'\alpha}_{j\lambda'_{2}\lambda_{2}}(p,m',m)\,a^{\dg b}_{j}(p)d^{\alpha'\dg}_{\lambda'_{2}}(m')d^{\alpha}_{\lambda_{2}}(m)
\;-\;\mathrm{h.c.}
\end{equation}
are absent from $\OM_{q\bar qg}$. They are $O(g)$, they are the quark analogues of the three-gluon
coefficient $C$, and without them the $O(g)$ diagonalization of the Hamiltonian cannot succeed:
the matrix elements of $H_{q\bar qg}$ connecting $\ket{q\bar q}$ and $\ket{q\bar q g}$ would remain
uncancelled.

If $\mf F$ and $\mf G$ are added to $\OM_{1}$, the results \eqref{eq:A4}, \eqref{eq:C0},
\eqref{eq:A2}, \eqref{eq:B2}, \eqref{eq:C1}, \eqref{eq:C2} and \eqref{eq:D2} are
\emph{unchanged}, while two relations acquire extra terms:
\begin{align}
\text{Eq.~\eqref{eq:D3}:}\quad&\text{add}\quad
-\ip{p}\Big[\mf F^{b\beta\beta'}_{j\lambda_{1}\lambda'_{1}}(p,l,l')\Big]^{\dg}
\mf G^{b\alpha\alpha'}_{j\lambda_{2}\lambda'_{2}}(p,m,m')+\big(\mf F\leftrightarrow\mf G\big),
\\
\text{Eq.~\eqref{eq:B4}:}\quad&\text{add}\quad
\ip{\mu}\mf F^{c\beta\mu}_{k\lambda_{1}\lambda_{\mu}}(q,l,l_{\mu})
\mf A^{b\alpha\mu}_{3j\lambda_{2}\lambda_{\mu}}(p,m,l_{\mu})
\nonumber\\&\hspace{2.1cm}
+\ip{\bar\mu}\mf G^{c\alpha\bar\mu}_{k\lambda_{2}\lambda_{\bar\mu}}(q,m,m_{\bar\mu})
\mf A^{b\bar\mu\beta}_{3j\lambda_{\bar\mu}\lambda_{1}}(p,m_{\bar\mu},l),
\end{align}
and new relations appear for the (also missing) $b^{\dg}b$ and $d^{\dg}d$ coefficients,
$\mathfrak P+\mathfrak P^{\dg}=-\int_{p}\int_{\kappa'}\mf F^{\dg}\mf F$ and likewise for
$\bar{\mathfrak P}$ with $\mf G$.
\end{warn}

%==================================================================
\section{Summary}
%==================================================================

\begin{center}
\renewcommand{\arraystretch}{1.3}
\begin{tabular}{lll}
\hline
coefficient & order & determined by\\
\hline
$\mf A_{3}$ & $g$ & LCPT ($g\to q\bar q$)\\
$\mf A_{4}$ & $g$ & \textbf{unitarity}, Eq.~\eqref{eq:A4}: $\;=-\mf A^{\dg}_{3}$\\
$\mf C_{0}$ & $g^{0}$ & \textbf{unitarity}, Eq.~\eqref{eq:C0}: $\;=1+O(g^{4})$\\
$\mf A_{1}$ & $g^{2}$ & LCPT (vacuum $\to q\bar q$)\\
$\mf A_{2}$ & $g^{2}$ & \textbf{unitarity}, Eq.~\eqref{eq:A2}: $\;\mf A_{1},\mf A_{3},A$\\
$\mf B_{1}$ & $g^{2}$ & LCPT ($gg\to q\bar q$)\\
$\mf B_{2}$ & $g^{2}$ & \textbf{unitarity}, Eq.~\eqref{eq:B2}: $\;\mf B_{1},\mf A_{3},A,C$\\
$\mf B_{3}$ & $g^{2}$ & LCPT ($g\to g\,q\bar q$)\\
$\mf B_{4}$ & $g^{2}$ & \textbf{unitarity}, Eq.~\eqref{eq:B4}: $\;\mf B_{3},\mf A_{3},A^{\dg},C^{\dg}$\\
$\mf C_{3}$ & $g^{2}$ & LCPT ($ggg\to g\,q\bar q$)\\
$\mf C_{1}$ & $g^{2}$ & \textbf{unitarity}, Eq.~\eqref{eq:C1}: $\;\mf C_{3},\mf A_{3},C$\\
$\mf C_{4}$ & $g^{2}$ & LCPT ($gg\to gg\,q\bar q$)\\
$\mf C_{2}$ & $g^{2}$ & \textbf{unitarity}, Eq.~\eqref{eq:C2}: $\;\mf C_{4},\mf A_{3},C^{\dg}$\\
$\mf D_{1}$ & $g^{2}$ & LCPT ($gg\to q\bar q\,q\bar q$)\\
$\mf D_{2}$ & $g^{2}$ & \textbf{unitarity}, Eq.~\eqref{eq:D2}: $\;\mf D_{1},\mf A_{3}$\\
$\mf D_{3}$ & $g^{2}$ & LCPT ($q\bar q\to q\bar q$), Hermitian part constrained by Eq.~\eqref{eq:D3}\\
\hline
\end{tabular}
\end{center}

\noindent
Structurally, every unitarity-fixed coefficient takes the form
\begin{equation}
[\text{unitarity-fixed}]=-\big[\text{its LCPT-fixed conjugate partner}\big]^{\dg}
+\big[\text{products of lower order coefficients}\big],
\end{equation}
where the second piece is present whenever the corresponding structure can be reached by a
product of two $O(g)$ transitions. Only $\mf A_{4}$ and $\mf C_{0}$ have no such product term at
the order considered.

\appendix
%==================================================================
\section{Explicit LCPT check of the $\mf A_{2}$ relation}
\label{app:check}
%==================================================================

Because Eq.~\eqref{eq:A2conj} is the relation most sensitive to the convention issue raised in
Section~\ref{sec:twosectors}, we verify it directly against the LCPT wave functions.

\subsection{Normalization rule}

If a term of $\OM$ has $n$ annihilation and $M$ creation operators with coefficient $K$, then
acting on the corresponding $n$-particle state,
\begin{equation}
\OM\ket{n}\supset(2\pi)^{-3(n+M)/2}\!\int\! d^{3}(\text{$M$ momenta})\;K\;\ket{M},
\quad\text{i.e.}\quad
K=(2\pi)^{\frac{3}{2}(n+M)}\big[\text{LCWF coeff.}\big].
\label{eq:rule}
\end{equation}
(This reproduces the factors $(2\pi)^{3/2}$ for $A$, $(2\pi)^{3}$ for $B_{2},B_{3}$,
$(2\pi)^{9/2}$ for $C$ and $(2\pi)^{6}$ for $D_{1}$ quoted in the pure gauge paper.) Hence
$\mf A_{1}=(2\pi)^{3}\Psi$ and $\mf A_{3}=(2\pi)^{9/2}\Psi$ for the respective wave function
coefficients.

\subsection{Ingredients}

With $\eps_{p}=\bm p^{2}/2p^{+}$, $k=l+m$, $\eps_{q\bar q}=\eps_{l}+\eps_{m}$, and
$\Phi^{i}_{\lambda_1\lambda_2}(k;l,m)=\chi^{\dg}_{\lambda_1}\Gamma^{i}(k;l,m)\chi_{\lambda_2}$,
the matrix elements are
\begin{align}
\bra{g^{a}_{i}(k)}H_{g}\vac&=\frac{g\,k^{i}\rho^{a}(-k)}{4\pi^{3/2}|k^{+}|^{3/2}},
&
\bra{q\bar q}H_{q\bar qg}\ket{g^{a}_{i}(k)}&=\frac{g\,t^{a}_{\alpha\beta}}{8\pi^{3/2}\sqrt{k^{+}}}
\dthree(k-l-m)\Phi^{i}_{\lambda_1\lambda_2},
\nonumber\\
\bra{q\bar q}H_{\rm inst}\vac&=\frac{g^{2}t^{a}_{\alpha\beta}\rho^{a}(-k)\delta_{\lambda_1\lambda_2}}
{8\pi^{3}(k^{+})^{2}}, &&
\end{align}
the last being the instantaneous $\rho\cdot\mathcal J_{q}$ exchange. Applying
Eq.~\eqref{eq:rule},
\begin{equation}
A^{b}_{j}(p)=-\frac{\sqrt2\,g\,p^{j}\rho^{b}(-p)}{\sqrt{p^{+}}\,p^{2}},
\qquad
\mf A_{3}=\frac{(2\pi)^{3}g\,t^{a}_{\alpha\beta}}{2\sqrt2\sqrt{k^{+}}}
\frac{\dthree(k-l-m)\,\Phi^{i}_{\lambda_1\lambda_2}(k;l,m)}{\eps_{k}-\eps_{l}-\eps_{m}} .
\end{equation}

\subsection{The two sides}

For the vacuum incoming state, second order perturbation theory ($\rho$ emits, gluon splits) plus
the first order instantaneous term give
\begin{equation}
\mf A_{1}=\frac{g^{2}t^{a}\rho^{a}(-k)\,k^{i}\Phi^{i}}{2k^{+}k^{2}\eps_{q\bar q}}
-\frac{g^{2}t^{a}\rho^{a}(-k)\,\delta_{\lambda_1\lambda_2}}{(k^{+})^{2}\eps_{q\bar q}} .
\end{equation}
For the $q\bar q$ incoming state the energy denominators are $(\eps_{q\bar q}-0)$ and
$(\eps_{q\bar q}-\eps_{k})$, giving
\begin{equation}
\mf A_{2}^{\dg}=\frac{g^{2}t^{a}\rho^{a}(-k)\,k^{i}\Phi^{i}}
{4(k^{+})^{2}\eps_{q\bar q}(\eps_{q\bar q}-\eps_{k})}
+\frac{g^{2}t^{a}\rho^{a}(-k)\,\delta_{\lambda_1\lambda_2}}{(k^{+})^{2}\eps_{q\bar q}} .
\end{equation}

\paragraph{Instantaneous terms.} They enter $\mf A_{1}$ with a denominator $(0-\eps_{q\bar q})$
and $\mf A_{2}$ with $(\eps_{q\bar q}-0)$, hence with opposite signs, and cancel exactly in the
sum $\mf A_{1}+\mf A_{2}^{\dg}$. This is required, since the right-hand side of
Eq.~\eqref{eq:A2conj} is built from $O(g)$ coefficients only and contains no instantaneous piece.

\paragraph{Remaining terms.} Using $4(k^{+})^{2}\eps_{k}=2k^{+}k^{2}$,
\begin{equation}
\frac{1}{2k^{+}k^{2}\eps_{q\bar q}}+\frac{1}{4(k^{+})^{2}\eps_{q\bar q}(\eps_{q\bar q}-\eps_{k})}
=\frac{1}{2k^{+}k^{2}\eps_{q\bar q}}\Big[1+\frac{\eps_{k}}{\eps_{q\bar q}-\eps_{k}}\Big]
=\frac{1}{2k^{+}k^{2}(\eps_{q\bar q}-\eps_{k})},
\end{equation}
so that
\begin{equation}
\mf A_{1}+\mf A_{2}^{\dg}
=\frac{g^{2}t^{a}\rho^{a}(-k)\,k^{i}\Phi^{i}_{\lambda_1\lambda_2}(k;l,m)}
{2k^{+}k^{2}\big(\eps_{q\bar q}-\eps_{k}\big)} .
\end{equation}
On the other side, the $\dthree$ in $\mf A_{3}$ collapses the integral,
\begin{equation}
\ip{p}\mf A_{3}(p,m,l)\,A^{b}_{j}(p)
=\frac{1}{(2\pi)^{3}}\cdot\frac{(2\pi)^{3}g\,t^{b}\Phi^{j}}{2\sqrt2\sqrt{k^{+}}(\eps_{k}-\eps_{q\bar q})}
\cdot\Big(-\frac{\sqrt2\,g\,k^{j}\rho^{b}(-k)}{\sqrt{k^{+}}k^{2}}\Big)
=\frac{g^{2}t^{b}\rho^{b}(-k)k^{j}\Phi^{j}}{2k^{+}k^{2}(\eps_{q\bar q}-\eps_{k})} .
\end{equation}

\begin{res}
The two sides agree identically. This confirms Eq.~\eqref{eq:A2conj}, the normalization rule
\eqref{eq:rule}, and --- crucially --- that the cross term with the pure gauge coefficient $A$
is genuinely present, so that $\mf A_{2}\neq-\mf A_{1}^{\dg}$.
\end{res}

\noindent
As a by-product one obtains a compact form for the pair creation coefficient itself. Using
$k^{i}\Gamma^{i}=4\eps_{k}-2\eps_{q\bar q}
-\frac{k^{+}}{l^{+}m^{+}}(\bm\sigma\!\cdot\!\bm l)(\bm\sigma\!\cdot\!\bm m)$, the instantaneous
term cancels the $\delta_{\lambda_1\lambda_2}/(k^{+})^{2}\eps_{q\bar q}$ piece exactly and
\begin{equation}
\mf A^{\alpha\beta}_{1\lambda_{2}\lambda_{1}}(m,l)
=-g^{2}t^{a}_{\alpha\beta}\rho^{a}(-k)\left[
\frac{\delta_{\lambda_1\lambda_2}}{k^{+}k^{2}}
+\frac{\chi^{\dg}_{\lambda_1}\big(\bm l\!\cdot\!\bm m+i(\bm l\times\bm m)\sigma^{3}\big)\chi_{\lambda_2}}
{2l^{+}m^{+}k^{2}\,\eps_{q\bar q}}\right],
\qquad \bm l\times\bm m\equiv\epsilon^{ij}l^{i}m^{j}.
\end{equation}

\end{document}
